\keepXColumns{}
\begin{tabularx}{\textwidth}{
  >{\raggedright\arraybackslash}p{1cm} 
  >{\raggedright\arraybackslash}p{2cm} 
  X
}
  \caption{Definisi Aktor Perangkat Lunak Audit Energi}\label{tbl:actor} \\
  \toprule
  \textbf{Kode} & \textbf{Aktor} & \textbf{Deskripsi} \\
  \midrule
  \endfirsthead

  \caption{Definisi Aktor Perangkat Lunak Audit Energi (lanjutan)} \\
  \toprule
  \textbf{Kode} & \textbf{Aktor} & \textbf{Deskripsi} \\
  \midrule
  \endhead

  \midrule
  \multicolumn{3}{l}{\textit{Bersambung ke halaman berikutnya}} \\
  %\bottomrule
  \endfoot

  \bottomrule
  \endlastfoot

  % === Table Data ===
  \hline
  A01 & Auditor Energi & 
  Pengguna utama sistem yang bertanggung jawab melakukan pengelolaan proyek audit, menginput dan memverifikasi data, serta meninjau hasil audit dalam bentuk \textit{dashboard} dan laporan. \\
  \hline
  A02 & Pengelola Gedung & 
  Pihak pengelola fasilitas gedung yang menggunakan sistem untuk melihat hasil audit energi, khususnya \textit{dashboard} dan laporan, sebagai dasar pengambilan keputusan operasional. \\
  \hline
\end{tabularx}